% Removed 29 March 2029
% From Introduction 
This paper formalizes this idea. We decompose the daily effective
temperature into a systematic component reflecting the local climate and
an idiosyncratic component reflecting microclimate variation. Bloom
occurs at a hitting time, [the date on which the cumulative sum first
crosses a plant-specific temperature threshold. This representation
places biological timing problems within the domain of stochastic
process theory, where the distribution of hitting times can be
characterized asymptotically, greatly simplifying the statistical laws
that describe the behavior of temperature-triggered biological events.]

A central implication of this framework is that synchrony depends on how
quickly temperatures are changing at the time the threshold is reached.
For example, many temperate climates exhibit two different regimes when
flowers first bloom in the spring. In the first regime---roughly the
winter months following the winter solstice (January--February)---average daily
temperatures are approximately constant over time, so cumulative
temperature grows at an approximately linear rate. In the second
regime---roughly the spring months around the equinox
(March--April)---average daily temperatures increase approximately
linearly over time, so cumulative temperatures grows quadratically.

% Maybe add below to caption of Figure 1? 
Figure 1 visualizes both regimes. The horizontal dashed line corresponds
with regime 1 (approximately constant average daily temperatures, linear
cumulative temperatures), while the upward-sloping dashed line
corresponds with regime 2 (approximately linear average daily
temperatures, quadratic cumulative temperatures). These dashed lines are
the ``signal'' that coordinates ecosystems. The ``noise'' is the
fluctuations in daily temperature (solid line) around the average
(dashed line), which creates variation in the event times. The extent to
which an ecosystem is coordianted depends on the strength of the signal
relative to the noise.

To see how, note that the cumulative temperature can be divided into the
cumulative signal (the sum of the averages) plus the cumulative noise
(the sum of the random deviations). When the distribution of the noise
is roughly the same in both regimes, the quadratic signal in regime 2
can much more quickly overwhelm the noise when compared to the linear
signal in regime 1. For this reason, nearby buds tend to cross the
threshold within a narrow time window when events occur in regime 2. In
contrast, when thresholds are reached in the regime 1, cumulative noise
can remain relatively influential so that small microclimate differences
translate into larger differences in threshold-crossing times.

%%

The remainder of the paper makes this intuition precise. In this
section, briefly review the literatures of stopped random walks and
ecology, and we state our two main findings in that context. In Section
2, we present the hitting time model and derive asymptotic
approximations under the two regimes. In Section 3, we present several
datasets confirming our findings empirically. In Section 4, we conclude
with a discussion. A sketch of the derivations, additional figures, and
simulations are contained in the Appendix A.1.

\subsection{1.1 Literature}\label{literature}

The observation that the timing of many biological processes are
determined by cumulative temperatures is largely credited to René
Réaumur, an early adopter of the thermometer in the mid-eighteenth
century. Adolphe Quetelet was the first to study the phenomenon
extensively, devoting an entire chapter to the law of the flowering
plants in his Letters addressed to HRH the Grand Duke of Saxe Coburg and
Gotha (Quetelet 1849, Letter Number 33), which Kotz and Johnson list
among one of eleven major breakthroughs in statistics prior to 1900.
Quetlet stressed the importance of idiosyncratic variation, which we
also emphasize in this paper. The law was well known throughout the
nineteenth century, influencing statisticians such as Francis Galton and
Florence Nightingale and waning with the rise of mathematical statistics
(Hacking 1990, 62). 

The law entered modern phenology through Quetelet's collaborator Charles
Fran\c{c}ois Antoine Morren, who coined the term phenology in an 1849
public lecture to describe the systematic study of periodically
recurring life-cycle events. [Today, the law is used largely unchanged
over two centuries. Scientists predict event times by tracking
accumulated heat above a base temperature.] This methodology is now
codified in the growing-degree-day (GDD) or temperature-sum model and
its variants.

[At the same time, the mathematical tools for analyzing stopped random
walks were being developed in probability and statistics.] Random-walk
models were formalized and popularized in the early twentieth century by
Pearson (1905) and Wald (1992). A central milestone for inference with
random indices is Anscombe's 1952 limit theory for sequential
estimation, which helped motivate a broad ``stopped random walk''
methodology for deriving asymptotics of hitting times and related
functionals. See for example Woodroofe (1982) and Gut (2012).

\subsection{1.2 Summary of Main
Findings}\label{summary-of-main-findings}


%% Stuff I tried harder to add but did not fit

% I like this but I don't think we can build a paper around microclimate variation given our current understanding of it (it has been studied and the evidence is not that convincing but it has not been studied well) -- it fits much better as a discussion point than something in the introduction 
% Careful: an idiosyncratic component reflecting microclimate variation
Yet coordination based on cumulative temperatures does not perfectly
synchronize springtime events. Bloom dates vary even among nearby buds
on the same tree or within the same grove. A leading source of this
variation is due to microclimates. The effective temperature experienced
by a bud depends on a litany of factors such as airflow, sun and water
exposure, and animal and human activity. These features introduce
idiosyncratic variation around a broader seasonal pattern so that the
cumulative temperature experienced by each bud is better viewed as a
random walk rather than a deterministic sum.


Warming advances threshold-triggered events because the cumulative total
reaches the target sooner. Typically, one would expect the rate of
advancement to diminish as temperatures warm. Indeed, when traveling a
fixed distance, one has to double the speed to half the travel time so
that the relationship between speed and time is a hyperbola.

The twist is that warming can change which aspects of the temperature
drives the timing of threshold-triggered events. If average
temperatures increase day over day as it does for springtime events in
temperate climates near the equinox, the temperature rate not the
baseline level of temperatures largely determined the bloom date. In
contrast, if average temperatures are relatively flat day over day as
they are in temperate climates closer to the solstice, the baseline
temperature determines the boom date. Thus, as warming moves springtime
events closer to the solstice, the rate at which their timing occurs can
increase, accelerating the advancement of spring

The second finding concerns spread rather than the average. Even if
warming makes the mean bloom date earlier, it does not necessarily make
timing more predictable. The date depends on both the average
temperature trend and idosynctatic variation induced in part by
microclimate factors such as shade, wind, soil moisture, urban heat, and
other factors that vary even within a single site. [When average
temperatures rise quickly and steadily, as when closer to the equinox,
those local fluctuations are relatively small compared with the strong
average temperature ``signal,'' so nearby buds tend to cross the
threshold within a narrow window. When average temperatures are
relatively constant, as when closer to the solstice, the seasonal signal
is weaker, and small local differences can have outsized effects. For
example, one bud may cross after an early warm spell while another must
wait for the next, widening the range of dates.]

Therefore, warming can both increase or decrease advancement and
variability depending on where in the seasonal cycle the threshold is
reached. In particular, if warming advances an event into an earlier
period with weaker accumulation, microclimate noise has more influence,
and synchrony can break down. [This regime shift produces a ``scattered
spring'' when climate change moves events earlier on average while
simultaneously reducing coordination across individuals and locations,
making spring less tightly synchronized even as it arrives sooner.]

